\documentclass[10pt,a4paper]{amsart}
\usepackage[margin=19mm]{geometry}
\usepackage{amsmath,amssymb,mathtools}
\usepackage[scr=rsfs,cal=euler]{mathalfa}
\usepackage{xfrac}
\usepackage[unicode,hidelinks,bookmarks=false]{hyperref}
\newcommand{\ie}{i.e.\ }
\newcommand{\cf}{cf.\ }
\newcommand{\resp}{respectively\ }
\newcommand{\Subsection}{\subsection}
\providecommand{\nobreakdash}{\nobreak\hbox{-}}
\setlength{\emergencystretch}{2em}
\multlinegap0pt
\usepackage{amssymb}
\usepackage{tikz}
\usepackage{xcolor}
\usepackage{float}
\usepackage{algorithm}\usepackage{algpseudocode}
\newtheorem{theorem}{Theorem}
\def\bs{\boldsymbol}
\def\th{\boldsymbol{\theta}}
\title{Empirical-Bayes Elastic-Net Computation for Exponential Random Graph Models}
\author{Dan Han, Vicki Modisette, Ting Li, Akidul Haque}
\date{Source: arXiv:2608.25280v1 (CC BY 4.0)}
\begin{document}
\maketitle

\noindent\emph{Source and licence.} Empirical-Bayes Elastic-Net Computation for Exponential Random Graph Models. Version arXiv:2608.25280v1 (CC BY 4.0), distributed by arXiv under the Creative Commons Attribution 4.0 International licence, http://creativecommons.org/licenses/by/4.0/, as recorded in the arXiv licence metadata for this exact version and shown on the abstract page https://arxiv.org/abs/2608.25280v1. Full licence notice: \texttt{licenses/arxiv-2608.25280-CC-BY-4.0.txt}.\par\smallskip

\noindent\textit{Editorial scope.} Counted unit: the article's theorem thm:full-likelihood-grouping-effect (source0.tex lines 1268--1297). The article's complete original proof of it is delivered with it (source0.tex lines 2538--2634, one \texttt{proof} environment of 97 lines, which the article places away from the statement). No mathematics is written, repaired or re-derived: the statement and the proof are the article's own lines, in the article's own notation, and no retained line is substituted or re-flowed.  No further source-local context is needed: every cross-reference of every retained line resolves inside the two delivered spans.  Nothing was omitted from any retained span, so the package carries no omission record.  No retained line contains a citation, so the wrapper carries no bibliography and no bibliography entry.  No retained line contains a figure environment, an image-inclusion command or drawing code, so no image is delivered and the package declares no asset dependency.  Editorial formatting only, in addition: an AMS \texttt{amsart} preamble that restates the article's own declarations and macros used by the retained lines, namely the environments \texttt{theorem} on separate counters, together with the standard packages those lines need.  The retained environments print in this document as Theorem 1 in order of appearance, whereas the article prints the same items with the numbering of its own full document.  The complete native source of the article as distributed in the arXiv e-print for this exact version is bundled unchanged as \texttt{sources/208564/source0.tex}.
\begin{theorem}[Full-likelihood MAP coefficient grouping]
\label{thm:full-likelihood-grouping-effect}
Suppose that \(\mathcal Y\) is finite,
\(\lambda_1\geq0\), and \(\lambda_2>0\). If
\[
\widehat\theta_j\neq0,\qquad
\widehat\theta_k\neq0,\qquad
\operatorname{sign}(\widehat\theta_j)
=
\operatorname{sign}(\widehat\theta_k),
\]
then
\[
2\lambda_2(\widehat\theta_j-\widehat\theta_k)
=
g_{jk}(\bs y)
-
\mathbb E_{\widehat\th}\{g_{jk}(\bs Y)\},
\]
and consequently
\[
|\widehat\theta_j-\widehat\theta_k|
\le
\frac{D_{jk}}{2\lambda_2}.
\]
In particular, if \(s_j-s_k\) is constant on \(\mathcal Y\), then
\[
\widehat\theta_j=\widehat\theta_k.
\]
\end{theorem}

\begin{proof}[Proof of
Theorem~\ref{thm:full-likelihood-grouping-effect}]
Consider the full-likelihood elastic-net objective
\[
\Phi(\th)
=
\log z(\th)-\th^\top s(\bs y)
+
\lambda_1\|\th\|_1
+
\lambda_2\|\th\|_2^2.
\]
Because \(\mathcal Y\) is finite, \(\log z(\th)\) is finite, smooth, and
convex. Since \(\lambda_2>0\), the objective \(\Phi\) is
\(2\lambda_2\)-strongly convex and coercive. It therefore has a unique
minimizer \(\widehat\th\).

Differentiating the log-partition function gives
\[
\frac{\partial}{\partial\theta_j}\log z(\th)
=
\mathbb E_{\th}\{s_j(\bs Y)\}.
\]
For a nonzero MAP coordinate \(\widehat\theta_j\), the KKT equation is
\[
\mathbb E_{\widehat\th}\{s_j(\bs Y)\}
-
s_j(\bs y)
+
\lambda_1\operatorname{sign}(\widehat\theta_j)
+
2\lambda_2\widehat\theta_j
=
0.
\]
Similarly, for coordinate \(k\),
\[
\mathbb E_{\widehat\th}\{s_k(\bs Y)\}
-
s_k(\bs y)
+
\lambda_1\operatorname{sign}(\widehat\theta_k)
+
2\lambda_2\widehat\theta_k
=
0.
\]
Because the two coordinates have the same sign, subtracting these equations
eliminates the \(L_1\) terms. It follows that
\[
2\lambda_2(\widehat\theta_j-\widehat\theta_k)
=
\{s_j(\bs y)-s_k(\bs y)\}
-
\mathbb E_{\widehat\th}
\{s_j(\bs Y)-s_k(\bs Y)\},
\]
or, equivalently,
\[
2\lambda_2(\widehat\theta_j-\widehat\theta_k)
=
g_{jk}(\bs y)
-
\mathbb E_{\widehat\th}\{g_{jk}(\bs Y)\}.
\]

Since \(\mathcal Y\) is finite, both
\(g_{jk}(\bs y)\) and
\(\mathbb E_{\widehat\th}\{g_{jk}(\bs Y)\}\) lie in the interval
\[
\left[
\min_{\bs u\in\mathcal Y}g_{jk}(\bs u),
\,
\max_{\bs u\in\mathcal Y}g_{jk}(\bs u)
\right].
\]
Therefore,
\[
\left|
g_{jk}(\bs y)
-
\mathbb E_{\widehat\th}\{g_{jk}(\bs Y)\}
\right|
\leq
D_{jk}.
\]
Dividing by \(2\lambda_2\) yields
\[
|\widehat\theta_j-\widehat\theta_k|
\leq
\frac{D_{jk}}{2\lambda_2}.
\]
If \(s_j-s_k\) is constant on \(\mathcal Y\), then \(D_{jk}=0\), and hence
\[
\widehat\theta_j=\widehat\theta_k.
\]
\end{proof}

\end{document}
